% Generated draft: requires mathematical and source review.
\documentclass[11pt]{article}
\usepackage[margin=25mm]{geometry}
\usepackage{fontspec}
\setmainfont{Latin Modern Roman}
\usepackage{amsmath,amssymb,mathtools,unicode-math}
\setmathfont{Latin Modern Math}
\usepackage{enumitem,needspace,xurl,hyperref}
\setlength{\parindent}{0pt}
\setlength{\parskip}{5pt}
\setlength{\emergencystretch}{3em}
\newcommand{\N}{\mathbb{N}}
\newcommand{\Z}{\mathbb{Z}}
\newcommand{\Q}{\mathbb{Q}}
\newcommand{\R}{\mathbb{R}}
\newcommand{\C}{\mathbb{C}}
\newcommand{\F}{\mathbb{F}}
\newcommand{\A}{\mathbb{A}}
\newcommand{\PP}{\mathbb{P}}
\newcommand{\E}{\mathbb{E}}
\newcommand{\eps}{\varepsilon}
\newcommand{\dd}{\,\mathrm{d}}
\begin{document}
\section*{UnsolvedMath ID: 30003192}
\textbf{Source problem code:} catalogue entry
\subsection{Definitions and mathematical statement}
Let \(K\) be an algebraically closed, complete, nontrivially valued field, let \(X\) be a projective variety over \(K\), and put \(n=\dim X\). The supplied verification note identifies the intended setting as non-Archimedean, although the displayed source statement does not explicitly impose that restriction on the valuation. Let \(\omega\) be a normalized positive form on \(X\), in the terminology of the source. The supplied record does not specify the analytic space, the class of forms, or the normalization condition.

For a line bundle \(M\) on \(X\), write \(\overline M=(M,\|\cdot\|)\) for a metrized line bundle and \(c_1(\overline M)\) for its curvature in the intended analytic theory. A difference-of-semipositive metrization means a metrization expressed through an isomorphism
\[
\overline M\simeq\overline L_1\otimes\overline L_2^{-1},
\]
where the metrized line bundles \(\overline L_1\) and \(\overline L_2\) are semipositive. The precise meaning and regularity requirements of semipositivity are not specified in the supplied record.

\textbf{Source question, with unresolved qualifications.} Does every line bundle \(M\) on \(X\) admit a difference-of-semipositive metrization \(\overline M\), unique up to multiplication of its metric by a positive real constant, such that \(c_1(\overline M)\) is \(\omega\)-bounded and
\[
c_1(\overline M)\,\omega^{n-1}=\lambda(M)\,\omega^n,
\qquad
\lambda(M)=c_1(M)\cdot[\omega]^n?
\]
Here the formula for \(\lambda(M)\) is reproduced exactly from the supplied statement. Under the usual interpretation of \([\omega]\) as a class of degree \((1,1)\), its displayed product with \(c_1(M)\) has total intersection degree \(n+1\) on an \(n\)-dimensional variety and does not give the stated scalar by ordinary top intersection. The record supplies no alternative convention resolving this issue, so the constant cannot be defined precisely from the available evidence.

The notions of \(\omega\)-boundedness and \(\omega\)-harmonicity, and the sense in which the displayed curvature products and equality are defined, also require the missing boundedness and Laplacian setup mentioned in the original statement. Moreover, the exponent \(n-1\) requires an interpretation or a positive-dimensional hypothesis when \(n=0\); no such hypothesis is supplied. These qualifications remain unresolved, and the displayed question is a statement-only reconstruction requiring mathematical review.
\subsection{Short English statement}
The source asks whether each line bundle admits a metrization, unique up to constant rescaling, that is a difference of semipositive metrics and has \(\omega\)-bounded, \(\omega\)-harmonic curvature satisfying the displayed equation. The analytic definitions, normalization, dimension restriction, and dimensionally problematic formula for \(\lambda(M)\) remain unresolved.
\subsection{Sources}
\begin{itemize}
\item Supplied local UnsolvedMath registry record, permanent ID \(30003192\), title \emph{Unique Harmonic Metrizations of Line Bundles}, exact source problem code \texttt{catalogue entry}, in the Oberwolfach Reports Open Problems collection. No individual proposer is named; the supplied proposed year is 2016.
\item Supplied report citation: \emph{Arithmetic Geometry (2017)}, Oberwolfach Reports, DOI 10.4171/owr/2016/38, pp.\ 2171--2224. \url{https://doi.org/10.4171/owr/2016/38}. The report attribution and its year wording are retained as supplied.
\item Supplied literature source: Yuan--Zhang, \emph{arithmetic Hodge index I}. \url{https://arxiv.org/abs/1304.3538}.
\item Supplied literature source: Yuan--Zhang, \emph{arithmetic Hodge index II}. \url{https://arxiv.org/abs/1304.3539}.
\item Supplied literature source: \emph{Oberwolfach Report 38}. \url{https://doi.org/10.4171/owr/2016/38}.
\item Literature assessment dated exactly \texttt{2026-08-21}: the supplied curation labels the problem partially solved and reports proofs for curves, model metrics, dual-complex cases, and polarized dynamical systems, while describing the general case as open. This is supplied dated curation, not an independently verified assessment; the record explicitly requests independent verification of time-sensitive status and cited claims.
\item The registry labels the statement \texttt{corrected\_verified} and reports that the official report was checked to restore the non-Archimedean projective setting, difference-of-semipositive metric, boundedness condition, normalization, and constant. This is a supplied source-verification claim, not verification performed for this draft.
\end{itemize}
\end{document}
